\section{System Architecture}
\label{sec:architecture}

\begin{figure}[t]
\centering
\begin{tikzpicture}[node distance=4mm and 6mm, every node/.style={font=\footnotesize},
  box/.style={draw, rounded corners=1pt, align=center, minimum width=34mm, minimum height=6mm},
  arr/.style={-{Stealth[length=2mm]}}]
\node[box] (sub) {Periodic evidence submission\\(alerts, cases, steps, escalations,\\dispositions, assets)};
\node[box, below=of sub] (val) {Validation and normalization\\(per-record provenance, all-or-nothing)};
\node[box, below=of val, minimum height=14mm] (ana) {Analytical workers (\V{CODE-workers})\\execution gap $\cdot$ negative space\\anomaly $\cdot$ peer deviation\\coverage $\cdot$ workflow $\cdot$ \ldots};
\node[box, below=of ana] (risk) {Correlation, risk (\V{CODE-dimensions} dimensions)\\and prioritization};
\node[box, right=12mm of sub] (find) {Evidence-cited findings\\and recommendations};
\node[box, below=of find] (rev) {Human supervisory review\\(review interrupt)};
\node[box, below=of rev] (dec) {Authorised decision};
\node[box, below=of dec] (trust) {TRUST-SAT finalization\\SHA3-256 canonical record,\\\V{CODE-signature} receipt, ledger};
\node[box, below=of trust] (ver) {Verification and audit};
\draw[arr] (sub) -- (val); \draw[arr] (val) -- (ana); \draw[arr] (ana) -- (risk);
\draw[arr] (risk.east) -| ++(4mm,0) |- (find.west);
\draw[arr] (find) -- (rev); \draw[arr] (rev) -- (dec); \draw[arr] (dec) -- (trust); \draw[arr] (trust) -- (ver);
\end{tikzpicture}
\caption{SAT-SA supervisory workflow as implemented. Analytics are deterministic;
orchestration (direct worker or LangGraph) sequences the stages and pauses for the
human decision.}
\label{fig:architecture}
\end{figure}

Figure~\ref{fig:architecture} shows the workflow; Table~\ref{tab:system} lists the
facts the experiments rely on, generated from the source code.

\paragraph{Submission and validation.}
An analyst of the monitored organisation creates a versioned submission for an
assessment period and uploads one file per evidence category through an HTTP API
(\V{CODE-api-routes} routes). Validation parses each record, resolves references
(an investigation step must reference an existing case), enforces chronology
(closure not before acknowledgement) and stores per-record source provenance
(artifact digest, record locator). Hosted validation is all-or-nothing: any
rejected record invalidates the version. Artifacts are stored in local storage or
an S3-compatible object store; metadata in SQLite or PostgreSQL.

\paragraph{Analysis.}
A separate worker process claims queued runs with a lease and executes
\V{CODE-workers} deterministic analytical workers over the version's canonical
records. Each emits findings with a rationale, statistic, threshold, confidence
vector and references to the source records it depends on; a signal finding
without evidence references is invalid and is withheld. No worker uses a language
model; LangGraph is used only for orchestration. The repository's capability
registry lists \V{CODE-registry} components (\V{CODE-registry-satsa} supervisory
components and \V{CODE-registry-mlops} retained machine-learning-operations agents);
the experiments exercise the \V{CODE-workers} default analytical workers and the
workflow around them, not every registered component.

\paragraph{Risk and prioritization.}
Findings are grouped into \V{CODE-dimensions} weighted dimensions (execution gap,
peer deviation, detection gap, negative space, anomaly, investigation quality,
escalation discipline). Each dimension's score saturates in the summed
confidence of its findings, and the total (capped at 100) orders entities for
supervisory attention. Findings receive recommendations.

\paragraph{Human review and orchestration.}
Runs can require review. In \emph{direct} mode the worker pauses the run after
analysis and recommendations; in \emph{graph} mode a LangGraph graph with
\V{CODE-graph-nodes} nodes (readiness, analysis, recommendations, human review,
trust boundary) checkpoints its state and interrupts at the review node~\citep{langgraph2026docs}.
Only a supervisor (or organisation administrator) can record the decision; the worker then
resumes.

\paragraph{TRUST-SAT finalization and verification.}
After the decision, the system builds a canonical supervisory document
(run, decision, finding, observation, risk and recommendation digests, submission
snapshot and source provenance), hashes it with SHA3-256, signs the digest with
\V{CODE-signature}, stores a receipt, and appends the event to a hash-chained
ledger. Verification rebuilds the document from live records and checks digest,
decision binding, signature and ledger chain (Section~\ref{sec:integrity}).

\paragraph{Deployment.}
The repository contains a container image, a compose topology (PostgreSQL,
SeaweedFS object storage, migration and key-initialisation jobs, API, worker) and a
CI job that would run it. All measurements in this paper were taken with SQLite and
local storage on one machine; the containerized topology, live PostgreSQL and live
S3 storage were not executed (Table~\ref{tab:deployment}).

% Generated by scripts/build_paper_assets.py from satsa-evidence-freeze-v2; do not edit.
% Sources: code facts and deployment record
\begin{table}[t]
\centering\small
\caption{SAT-SA components exercised by the experiments (code facts).}
\label{tab:system}
\begin{adjustbox}{max width=\linewidth}
\begin{tabular}{ll}
\toprule
Component & Implementation fact \\
\midrule
Analytical workers per run & \V{CODE-workers} deterministic workers \\
Risk model & \V{CODE-dimensions} weighted dimensions, saturating fusion of finding confidence \\
Orchestration & direct worker, or LangGraph graph with \V{CODE-graph-nodes} nodes and a review interrupt \\
Decision binding & \V{CODE-signature} signature over a SHA3-256 canonical digest; hash-chained ledger \\
Interfaces & \V{CODE-api-routes} HTTP API routes; separate analysis worker process \\
Persistence & SQLite (all measurements); PostgreSQL path implemented, not live-tested \\
\bottomrule
\end{tabular}
\end{adjustbox}
\end{table}

